\documentclass[11pt,a4paper]{article}
\usepackage[T1]{fontenc}\usepackage[utf8]{inputenc}\usepackage{mathptmx}\usepackage[margin=2.5cm]{geometry}\usepackage{parskip}
\pagestyle{empty}
\begin{document}
\noindent Achaire Zogo Abouma Zozime (corresponding author)\\ University of Yaound\'e I\\ Yaound\'e, Cameroon\\ achaire.zogo@facsciences-uy1.cm\\[1em]
\today\\[1em]
The Editor-in-Chief\\ \textit{FinTech} (MDPI)\\[1em]
Dear Editor,


On behalf of my co-author, Edit Kamel Djupsa, and myself, please consider our manuscript, \textit{``Measuring the Operational Resilience of FinTech Ecosystems in Sub-Saharan Africa: The AFRO-FINRES Framework and Evidence from Cameroon and the CEMAC Region''}, for publication as an original research Article in \textit{FinTech}.

Digital financial services, especially mobile money, have become critical infrastructure in Sub-Saharan Africa. However, there is no comparative and replicable measure of whether national FinTech ecosystems can withstand and recover from cyber incidents, provider failures and infrastructure outages. Our manuscript fills this gap. It proposes AFRO-FINRES, a six-dimension framework that combines prevention, detection, response, recovery and continuity, governance and third-party infrastructure redundancy. We operationalise the framework as a composite index built only from public data (ITU, World Bank and PeeringDB) for 48 countries.

The main findings are as follows. The index is internally consistent and robust to alternative weighting and aggregation choices. It is associated with account and mobile-money ownership after controlling for income. The CEMAC region lags the regional average by about 13 points, and detection capability is Cameroon's binding constraint. An illustrative stress-test module shows how supervisors can use the framework to rank resilience investments. All code and data are provided for full replication.

We believe the manuscript fits \textit{FinTech}, as it addresses the risks, resilience and regulation of financial technology. It will interest researchers, regulators and practitioners working on digital finance in emerging markets.

We confirm that this manuscript is original, has not been published elsewhere and is not under consideration by another journal. The authors declare no conflicts of interest. Potential reviewers can be suggested on request.


Yours sincerely,\\[2em]
Achaire Zogo Abouma Zozime
\end{document}
